\section*{第 \ref{ch:g11:reaction-energy} 章——反应的能量：燃烧与键能}

\begin{solution}{exo:g11:reaction-energy:1}
木材燃烧：放热的。冰熔化：吸热的（它从
外界吸收热量，尽管它不是化学反应）。冷敷袋：
吸热的。暖手宝：放热的。阳光下制造糖：
吸热的（能量来自光）。
\end{solution}

\begin{solution}{exo:g11:reaction-energy:2}
$Q = 3.0 \times 890.7 = \qty{2.7e3}{kJ}$，约 \qty{2.7}{MJ}。
\end{solution}

\begin{solution}{exo:g11:reaction-energy:3}
反应物。它们的能量差主要以
热量形式释放到外界。
\end{solution}

\begin{solution}{exo:g11:reaction-energy:4}
$n = 100 / 46.0 = \qty{2.17}{mol}$；$Q = 2.17 \times 1367.6 =
\qty{2.97e3}{kJ}$，约 \qty{3.0}{MJ}。
\end{solution}

\begin{solution}{exo:g11:reaction-energy:5}
断开 1 摩尔这类键所需的能量（分子和原子
均为气态）。共用电子对将两个原子维系在一起，将它们
拉开需要克服这种吸引作用，因此需要能量。
\end{solution}

\begin{solution}{exo:g11:reaction-energy:6}
断键：$2 \times 436 + 498 = \qty{1370}{kJ}$；成键：
$4 \times 464 = \qty{1856}{kJ}$；$E_r \approx 1370 - 1856 =
\qty{-486}{kJ}$，放热的。
\end{solution}

\begin{solution}{exo:g11:reaction-energy:7}
乙醇 \ce{CH3-CH2-O-H} 含有 5 个 \ce{C-H}、1 个 \ce{C-C}、1 个 \ce{C-O} 和 1
个 \ce{O-H}。断键：$5 \times 415 + 345 + 350 + 464 + 3 \times 498 =
\qty{4728}{kJ}$。成键：$4 \times 741 + 6 \times 464 = \qty{5748}{kJ}$。
$E_r \approx \qty{-1020}{kJ}$，比
实测值 \qty{1367.6}{kJ} 的绝对值小约四分之一。
\end{solution}

\begin{solution}{exo:g11:reaction-energy:8}
丙烷：$2219.2 / 0.0440 = \qty{5.04e4}{kJ/kg} = \qty{50.4}{MJ/kg}$。
甲烷：$890.7 / 0.0160 = \qty{55.7}{MJ/kg}$。两者均与图中所示一致。
\end{solution}

\begin{solution}{exo:g11:reaction-energy:9}
$Q = 1500 \times 4.18 \times 85 = \qty{5.33e5}{J} = \qty{533}{kJ}$；
$n = 533 / 890.7 = \qty{0.598}{mol}$，即 $0.598 \times 16.0 =
\qty{9.6}{g}$ 的甲烷；若一半能量损失，则需 \qty{19}{g}。
\end{solution}

\begin{solution}{exo:g11:reaction-energy:10}
氢气 > 甲烷 > 丙烷 > 辛烷 > 乙醇。
$48.0 / 142.9 = \qty{0.34}{kg}$ 的氢气。
\end{solution}

\begin{solution}{exo:g11:reaction-energy:11}
断键：$946 + 3 \times 436 = \qty{2254}{kJ}$；成键：
$6 \times 390 = \qty{2340}{kJ}$；$E_r \approx \qty{-86}{kJ}$，略
放热。
\end{solution}

\begin{solution}{exo:g11:reaction-energy:12}
水：$250 \times 4.18 \times 20 = \qty{2.09e4}{J} = \qty{20.9}{kJ}$。
乙醇：$1.20 / 46.0 = \qty{0.0261}{mol}$，可释放
$0.0261 \times 1367.6 = \qty{35.7}{kJ}$。效率
$20.9 / 35.7 \approx \qty{59}{\percent}$。损失包括：热量被
空气带走，热量使金属罐和支架升温，燃烧不完全（产生炭黑），
以及部分乙醇未燃烧即挥发。
\end{solution}

\begin{solution}{exo:g11:reaction-energy:13}
\ce{2C8H18 + 25O2 -> 16CO2 + 18H2O}：生成 8 摩尔二氧化碳，
质量为 \qty{352}{g}，对应 \qty{5470}{kJ}，故 $352 / 5.470 = \qty{64.4}{g}$ 每
兆焦。\ce{C3H8 + 5O2 -> 3CO2 + 4H2O}：\qty{132}{g} 对应
\qty{2219.2}{kJ}，故每兆焦为 \qty{59.5}{g}。甲烷
每兆焦仅 \qty{49.4}{g}，释放的二氧化碳最少。
\end{solution}

\begin{solution}{exo:g11:reaction-energy:14}
表中给出的是平均键能，而二氧化碳中的 \ce{C=O} 键
比平均值更强；实测值针对的是
液态水，其液化释放的能量比气态
反应更多；并且该方法将每个键视为独立于其
相邻键。这一估算仍然有用：它给出了正确的符号
（放热）和正确的数量级，无需任何测量。
\end{solution}

\begin{solution}{exo:g11:reaction-energy:15}
断键：$4 \times 464 = \qty{1856}{kJ}$；成键：$2 \times 436 + 498 =
\qty{1370}{kJ}$；$E_r \approx \qty{+486}{kJ}$，吸热，因此需要
提供能量。之后燃烧氢气可将这部分能量重新释放出来：
氢气储存的是在别处产生的能量，它本身并不产生能量。
\end{solution}

\begin{solution}{pb:g11:reaction-energy:1}
\textbf{1.} \ce{CH4 + 2O2 -> CO2 + 2H2O}。

\textbf{2.} \ce{C2H6O + 3O2 -> 2CO2 + 3H2O}。

\textbf{3.} \ce{2C8H18 + 25O2 -> 16CO2 + 18H2O}。

\textbf{4.} \ce{2H2 + O2 -> 2H2O}。

\textbf{5.} 氢气。

\textbf{6.} \qty{16.0}{g/mol}、\qty{46.0}{g/mol}、\qty{114.0}{g/mol}、
\qty{2.0}{g/mol}。

\textbf{7.} 单位依次为 \si{kJ} 每 \si{kg}，再为 \si{MJ/kg}：甲烷
$890.7 / 0.0160$，\qty{55.7}{MJ/kg}；乙醇 $1367.6 / 0.0460$，
\qty{29.7}{MJ/kg}；辛烷 $5470 / 0.1140$，\qty{48.0}{MJ/kg}；氢气
$285.8 / 0.0020$，\qty{142.9}{MJ/kg}。

\textbf{8.} 氢气 > 甲烷 > 辛烷 > 乙醇。

\textbf{9.} 氢气是很轻的气体：一千克氢气在常压下
占据很大的体积，因此必须压入厚重的
高压储罐中。

\textbf{10.} 断裂的键：甲烷 4 个 \ce{C-H}；氧气 2 个 \ce{O=O}。
生成的键：二氧化碳 2 个 \ce{C=O}；水 $2 \times 2 = 4$ 个 \ce{O-H}。

\textbf{11.} $4 \times 415 + 2 \times 498 = \qty{2656}{kJ}$。

\textbf{12.} $2 \times 741 + 4 \times 464 = \qty{3338}{kJ}$。

\textbf{13.} $E_r \approx 2656 - 3338 = \qty{-682}{kJ}$，
与实测值 \qty{-890.7}{kJ} 相比，绝对值偏小约 \qty{23}{\percent}。

\textbf{14.} 估算所用的是平均键能（二氧化碳中的 \ce{C=O} 比平均值
更强）；实测中的水为液态，
估算中的物质为气态。

\textbf{15.} $1000 / 890.7 = \qty{1.123}{mol}$。

\textbf{16.} $1.123 \times 44.0 = \qty{49.4}{g}$。

\textbf{17.} 乙醇：$1000 / 1367.6 = \qty{0.731}{mol}$，可得
\qty{1.462}{mol} 二氧化碳，质量为 \qty{64.3}{g}。辛烷：
$1000 / 5470 = \qty{0.183}{mol}$，可得 \qty{1.463}{mol}，\qty{64.4}{g}。

\textbf{18.} 甲烷：每个碳原子所含氢原子最多（4 个，
辛烷为 2.25 个），且燃烧氢气可放能，
而不产生二氧化碳。

\textbf{19.} 关于氢气的制法：用可再生能源发的电
电解水制得，几乎不释放二氧化碳；由甲烷制得，
则二氧化碳在工厂排放，而非经
废气排出。

\textbf{20.} 甲烷每兆焦约放出
\qty{49}{g} 二氧化碳。
\end{solution}
